% ==============================================================================
% sections/data.tex
% Created: 2026-05-08
% ==============================================================================

\subsection{The Quarterly Labour Force Survey}

The Quarterly Labour Force Survey (QLFS) is a household-based sample survey
conducted by Statistics South Africa, introduced in 2008 to replace the
biannual Labour Force Survey. The survey is representative of the civilian
non-institutionalised population of South Africa and collects information on
employment status, earnings, industry, occupation, and individual and household
characteristics. It has a rotating panel structure that allows a subset of
individuals to be followed longitudinally.

We use the QLFS panel for the period 2010Q1 to 2019Q4. Individuals are matched
across consecutive waves using their household and person identifiers.
Statistics South Africa undertakes some consistency cleaning of the data prior
to release \citep{statssa2013}, but there is well-documented evidence of
remaining measurement problems in the matched panel
\citep{burger2016,lechtenfeld2014}. Our sample is restricted to individuals
aged 18 to 55 who have non-missing employment status in all available waves.
We impose no restriction on gender, race, or other demographics, to ensure
the results are representative of the working-age population as a whole.

For the three-wave baseline analysis, we retain all individuals observed in
three consecutive quarterly waves, yielding a sample of approximately 405,000
person-wave triads. For the four-wave AR(2) analysis in
Section~\ref{sec:alternatives}, we additionally require a valid fourth wave,
yielding approximately 317,000 observations. Survey weights are taken as the
geometric mean of the three (or four) wave-specific weights to produce a
balanced cross-wave weight \citep{kalbfleisch1985}.

\subsection{The empirical puzzle}\label{sec:data:puzzle}

Table~\ref{tab:descriptive} presents summary statistics for the three-wave
sample. Panel A shows that the sample spans the full working-age range, with
a mean age of approximately 34 years and a mean education level of roughly
10 years. The employment rate is stable across waves at approximately 48\%,
consistent with South Africa's persistently high non-employment rate.

Panel B documents the transition rate puzzle that motivates our approach.
The observed three-month entry and exit rates---the probability of moving from
non-employment to employment, and vice versa, between two consecutive quarterly
waves---are approximately 8\% in both directions. Under the first-order Markov
assumption, these three-month rates imply six-month rates of approximately
$1 - (1-0.08)^2 \approx 15\%$. But the six-month rates observed directly in
the four-wave data---the probability of transitioning over a six-month window,
using only the first and third waves---are only around 10\%. The implied
six-month rate exceeds the observed six-month rate by a factor of approximately
1.5. This pattern---where transition rates estimated from high-frequency panels
imply far more churn than lower-frequency estimates---is exactly the signature of
transient misclassification error \citep{meyer1988,porteba1986}. Under the
first-order Markov assumption on the \emph{true} process, a transient error at
wave $t$ creates a spurious entry and a spurious exit, causing the high-frequency
(3-month) estimate to be inflated while the low-frequency (6-month) estimate is
less affected because the error has reversed by the time the second wave is
observed.

This divergence is also consistent with the broader pattern in the South African
labour market literature, where high-frequency panels consistently yield higher
estimated transition rates than lower-frequency panels
\citep{banerjee2008,cichello2014}. \citet{burger2016} shows that the Markov
assumption applied directly to observed employment in South Africa is rejected
by the data, with classical transition estimates failing to fit the observed
three-wave moments. Our structural estimator attributes this violation to
misclassification error, while simultaneously testing the alternative explanations
of worker heterogeneity and duration dependence.

\subsection{Evidence of measurement error}

Panel C of Table~\ref{tab:descriptive} provides direct evidence that the QLFS
panel contains substantial measurement error, using two variables---age and
education---for which errors are more easily detected than in employment status.
An age inconsistency is defined as a wave-to-wave change in reported age that
is not in $\{0, 1\}$ years; since quarterly surveys are conducted 0.25 years apart,
age should remain constant or increase by one because responses are rounded to the nearest
integer year. An education
inconsistency is a wave-to-wave change that is negative or exceeds one year,
since education should be non-decreasing and cannot advance by more than one
level per quarter. These inconsistencies represent clear measurement error:
biological age cannot decrease, and educational attainment cannot be removed.

We find that approximately 11.6\% of individuals in our sample have at least one
age inconsistency across any pair of consecutive waves, and 26.8\% have at least
one education inconsistency. Overall, nearly one in three panel observations
(31.2\%) contains an inconsistent response in at least one of these two verifiable
variables. These figures, which are, if anything, an underestimate of overall
measurement error since they exclude undetectable errors in age (e.g., a plausible
but false age of 35 reported as 34), suggest that measurement error is widespread
in the QLFS panel. Crucially, we exploit the link between response unreliability
and employment misclassification in Section~\ref{sec:results}, where we show that
survey respondents who produce age or education inconsistencies also exhibit
employment dynamics that are consistent with significantly higher misclassification
rates.

\subsubsection{Tenure and duration data}\label{sec:data:tenure}

In addition to employment status, the QLFS collects information on job tenure
(how long the employed person has been with their current employer) and
non-employment duration (how long since the non-employed person last worked).
These duration variables are used in the tenure contamination model described
in Section~\ref{subsec:tenure}. They also provide direct
evidence of employment misclassification. Table~\ref{tab:tenure_diagnostics}
summarises three diagnostic moments from the linked panel that make this
evidence visible in the raw data.

Panel A begins with observed tenure, $g_t$ in period $t$. Job tenure evolves
deterministically: a worker who remains employed with the
same employer between two consecutive quarterly waves should report a tenure
increase of exactly three months (i.e. 0.25 years). Among all observed employed--employed
continuation pairs in the sample, approximately 65\% report precisely this
expected increment. The remaining 35\% report tenure changes that deviate from
the three-month clock---including negative changes, zero changes, and changes
that exceed a year---in a pattern consistent with a contamination process as in 
\citet{goldstein1982contamination} and \citet{horowitz1995identification}
rather than the classic "errors-in-variables" approach that 
would add, for example, a mean-zero Gaussian measurement noise. 
This two-component structure, in which
most reports are exact and a non-trivial minority are drawn from a dispersed
distribution, motivates the point-mass--plus--contamination measurement model
that we formalise in Section~\ref{subsec:tenure}.

Panel B shows that tenure data also reveal direct evidence that observed
employment status changes are sometimes spurious. Consider the subset of
individuals observed with the
sequence employed--non-employed--employed--- $(s_1, s_2, s_3) = (1,0,1)$--- who
appear to have lost a job and then found a new one. If the job loss at wave~2
was genuine and the worker was re-employed at wave~3, tenure at wave~3 should
be short---no more than three months. Instead, for nearly 30\% of these
individuals, tenure at wave~3 is approximately equal to tenure at wave~1 plus
six months, exactly the progression expected if the worker was continuously
employed throughout and the non-employment report at wave~2 was a
misclassification. The tenure flanks are thus informative about which
$(1,0,1)$ sequences reflect genuine transitions and which reflect classification
error.

Panel C reports a complementary pattern among within-panel ``fresh
starts''---individuals
observed as non-employed at wave~$t-1$ and employed at wave~$t$. If these
transitions were genuine, the worker would have started a new job within the
past three months, and reported tenure should be short. Yet approximately 37\%
of these apparent new entrants report tenure exceeding twelve months. A genuine
job entrant cannot have more than one year of tenure with the current employer;
these observations are therefore direct evidence of either misclassification of
employment status at wave~$t-1$ (the worker was actually employed all along) or
measurement error in reported tenure. Either way, they confirm that observed
employment transitions contain a substantial spurious component that is
corroborated by a variable---tenure---that the employment status question does
not directly reference.
